\NormalFont\ShortTitle{Marcos}
{\MT Evangelho Segundo Marcos

\par }\ChapOne{1}{\PP \VerseOne{1}Princípio do Evangelho de Jesus Cristo, Filho de Deus.
\VS{2}Como está escrito nos profetas: Eis que eu envio o meu mensageiro diante de tua face, que preparará o teu caminho diante de ti .
{\RQ{Malaquias 3:1
}}
\VS{3}Voz do que clama no deserto: Preparai o caminho do Senhor, endireitai suas veredas.
{\RQ{Isaías 40:3
}}
\VS{4}João veio a batizar no deserto, e a pregar o batismo de arrependimento para perdão dos pecados.
\VS{5}E toda a província da Judeia e os de Jerusalém saíam até ele; e eram todos batizados por ele no rio Jordão, confessando os seus pecados.
\VS{6}João se vestia de pelos de camelo, e um cinto de couro em sua cintura; e comia gafanhotos e mel do campo.
\VS{7}Ele pregava assim: Após mim vem aquele que é mais forte do que eu. A ele não sou digno de me abaixar para desatar a tira das suas sandálias.
\VS{8}De fato eu tenho vos batizado com água, porém ele vos batizará com Espírito Santo.
\VS{9}E aconteceu que, naqueles dias, Jesus de Nazaré da Galileia veio, e foi batizado por João no Jordão.
\VS{10}E assim que saiu da água, viu os céus se abrirem, e o Espírito, que como pomba descia sobre ele.
\VS{11}E veio uma voz dos céus: Tu és meu Filho amado, em quem me agrado.
\VS{12}E logo o Espírito o impeliu ao deserto.
\VS{13}Ele esteve ali no deserto quarenta dias, tentado por Satanás. Ele estava com os animais selvagens, e os anjos o serviam.
\VS{14}Depois que João foi preso, Jesus veio para a Galileia, pregando o Evangelho do Reino de Deus,
\VS{15}e dizendo: O tempo se cumpriu, e o Reino de Deus está perto; arrependei-vos, e crede no Evangelho.
\VS{16}E enquanto andava junto ao mar da Galileia, ele viu Simão e seu irmão André, que lançavam uma rede ao mar, porque eram pescadores;
\VS{17}Jesus lhes disse: Vinde após mim, e farei serdes pescadores de gente.
\VS{18}Então logo deixaram suas redes, e o seguiram.
\VS{19}E passando dali um pouco mais adiante, viu Tiago
{\ADD{filho}} de Zebedeu, e seu irmão João, que estavam no barco, consertando as redes.
\VS{20}E logo os chamou; então eles deixaram o seu pai Zebedeu no barco com os empregados, foram após ele.
\VS{21}Eles entraram em Cafarnaum; e assim que chegou o sábado,
{\ADD{Jesus}} entrou na sinagoga e começou a ensinar.
\VS{22}E ficavam admirados com o seu ensinamento, pois, diferentemente dos escribas, ele os ensinava como quem temautoridade.
\VS{23}E havia na sinagoga deles um homem com um espírito imundo, que gritou,
\VS{24}dizendo: Ah, que temos contigo, Jesus Nazareno? Vieste para nos destruir? Bem sei quem és: o Santo de Deus.
\VS{25}Jesus o repreendeu, dizendo: Cala-te, e sai dele.
\VS{26}E o espírito imundo, provocando convulsão nele, e gritando em alta voz, saiu dele.
\VS{27}Assim todos ficaram admirados, e perguntavam entre si: Que é isto? Que novo ensinamento é este que, com autoridade, ordena até aos espíritos imundos, e eles lhe obedecem?
\VS{28}E logo sua fama se espalhou por toda a região da Galileia.
\VS{29}Logo depois de saírem da sinagoga, vieram à casa de Simão e de André, com Tiago e João.
\VS{30}A sogra de Simão estava deitada com febre, e logo falaram dela
{\ADD{a Jesus}} .
\VS{31}Então ele aproximou-se dela, tomou-a pela mão, e a levantou; logo a febre a deixou, e ela começou a servi-los.
\VS{32}Ao entardecer, quando o sol já se punha, trouxeram-lhe todos os doentes e endemoninhados;
\VS{33}e toda a cidade se juntou à porta.
\VS{34}Ele curou muitos que se achavam mal de diversas enfermidades, e expulsou muitos demônios. Ele não deixava os demônios falarem, porque o conheciam.
\VS{35}De madrugada, ainda escuro, ele se levantou para sair, foi a um lugar deserto, e ali esteve a orar.
\VS{36}Simão e os que estavam com ele o seguiram.
\VS{37}Quando o acharam, disseram-lhe: Todos estão te procurando.
\VS{38}{\ADD{Jesus}} lhes respondeu: Vamos para as aldeias vizinhas, para que eu também pregue ali, pois vim para isso.
\VS{39}Ele pregava em suas sinagogas por toda a Galileia, e expulsava os demônios.
\VS{40}Um leproso aproximou-se dele, rogando-lhe, pondo-se de joelhos diante dele, e dizendo-lhe: Se quiseres, tu podes limpar-me.
\VS{41}E Jesus , movido de compaixão, estendeu a mão, tocou-o, e disse-lhe: Quero; sê limpo.
\VS{42}Quando disse isso, logo a lepra saiu dele, e ficou limpo.
\VS{43}{\ADD{Jesus}} advertiu-o, e logo o despediu,
\VS{44}dizendo-lhe: Cuidado, não digas nada a ninguém. Mas vai, mostra-te ao Sacerdote, e oferece por teres ficado limpo o que Moisés mandou, para lhes servir de testemunho.
\VS{45}Porém, quando ele saiu, começou a anunciar muitas coisas, e a divulgar a notícia, de maneira que
{\ADD{Jesus}} já não podia entrar publicamente na cidade; em vez disso, ficava do lado de fora em lugares desertos, e
{\ADD{pessoas}} de todas as partes vinham até ele.

\par }\Chap{2}{\PP \VerseOne{1}Dias depois,
{\ADD{Jesus}} entrou outra vez em Cafarnaum, e ouviu-se que estava em casa.
\VS{2}Logo juntaram-se tantos, que nem mesmo perto da porta cabiam; e ele lhes falava a palavra.
\VS{3}E vieram a ele uns que traziam um paralítico carregado por quatro.
\VS{4}Como não podiam se aproximar dele por causa da multidão, descobriram o telhado onde ele estava, fizeram um buraco, e baixaram
{\ADD{por ele}} o leito em que jazia o paralítico.
\VS{5}Quando Jesus viu a fé deles, disse ao paralítico: Filho, os teus pecados te são perdoados.
\VS{6}E estavam ali sentados alguns escribas, que pensavam
\FTNT{*}{{\FR 2:6 }
Lit. “indagavam”, “questionavam”
} em seus corações:
\VS{7}Por que este
{\ADD{homem}} fala essas blasfêmias? Quem pode perdoar pecados, a não ser somente Deus?
\VS{8}Imediatamente Jesus percebeu em seu espírito que assim pensavam em si mesmos. Então perguntou-lhes: Por que pensais assim
\FTNT{†}{{\FR 2:8 }
Lit. “essas coisas”
} em vossos corações?
\VS{9}O que é mais fácil? Dizer ao paralítico: “Os teus pecados estão perdoados”, ou dizer, “Levanta-te, toma o teu leito, e anda”?
\VS{10}Mas para que saibais que o Filho do homem tem poder na terra para perdoar pecados,(disse ao paralítico):
\VS{11}A ti eu digo: levanta-te, toma o teu leito, e vai para a tua casa.
\VS{12}E logo ele se levantou, tomou o leito, e saiu na presença de todos, de tal maneira, que todos ficaram admirados, e glorificaram a Deus, dizendo: Nunca vimos algo assim.
\VS{13}{\ADD{Jesus}} voltou a sair para o mar; toda a multidão veio até ele, e ele os ensinava.
\VS{14}E enquanto passava, ele viu Levi,
{\ADD{filho}} de Alfeu, sentado no posto de coleta de impostos, e disse-lhe: Segue-me.Então
{\ADD{Levi}} se levantou e o seguiu.
\VS{15}E aconteceu que enquanto estava sentado à mesa na casa dele, muitos publicanos e pecadores também estavam assentados à mesa com Jesus e os seus discípulos; porque eram muitos, e o haviam seguido.
\VS{16}Quando os escribas e os fariseus o viram comer com os publicanos e pecadores, disseram a seus discípulos: Por que é que ele come e bebe com os publicanos e os pecadores?
\VS{17}Jesus ouviu e lhes respondeu: Os que têm saúde não precisam de médico, mas sim os que estão doentes. Eu não vim para chamar os justos, mas sim, os pecadores, ao arrependimento.
\VS{18}Os discípulos de João e os dos fariseus estavam jejuando; então vieram lhe perguntar: Por que os discípulos de João e os dos fariseus jejuam, e os teus discípulos não jejuam?
\VS{19}Jesus lhes respondeu: Podem os convidados
\FTNT{‡}{{\FR 2:19 }
lit. filhos
} do casamento jejuar enquanto o noivo estiver com eles? Enquanto tiverem o noivo consigo, eles não podem jejuar.
\VS{20}Mas dias virão, quando o noivo lhes for tirado; e então naqueles dias jejuarão.
\VS{21}Ninguém costura remendo de pano novo em roupa velha; senão o remendo novo rompe o velho, e se faz pior rasgo.
\VS{22}E ninguém põe vinho novo em odres velhos; senão o vinho novo rompe os odres, derrama-se o vinho, e os odres se danificam; mas o vinho novo deve ser posto em odres novos.
\VS{23}E aconteceu que, enquanto
{\ADD{Jesus}} passava pelas plantações no sábado, os seus discípulos, andando, começaram a arrancar espigas.
\VS{24}Os fariseus lhe disseram: Olha! Por que estão fazendo o que não é lícito no sábado?
\VS{25}E ele lhes disse: Nunca lestes o que fez Davi, quando teve necessidade e fome, ele e os que com ele estavam?
\VS{26}Como ele entrou na Casa de Deus, quando Abiatar era sumo sacerdote, e comeu os pães da proposição (dos quais não é lícito comer, a não ser aos sacerdotes), e também deu aos que com ele estavam?
\VS{27}Disse-lhes mais: O sábado foi feito por causa do ser humano, não o ser humano por causa do sábado.
\VS{28}Por isso o Filho do homem é Senhor até do sábado.

\par }\Chap{3}{\PP \VerseOne{1}{\ADD{Jesus}} entrou outra vez na sinagoga; e estava ali um homem que tinha uma mão definhada.
\VS{2}E prestavam atenção nele, se o curaria no sábado, para o acusarem.
\VS{3}E
{\ADD{Jesus}} disse ao homem que tinha a mão definhada: Levanta-te, e vem para o meio.
\VS{4}E disse-lhes:É lícito no sábado fazer o bem, ou o mal? Salvar uma pessoa, ou matá-la? E mantiveram-se calados.
\VS{5}E olhando ao redor para eles com indignação, sentindo pena da dureza dos seus corações, disse ao homem: Estende a tua mão. E ele estendeu; e sua mão foi restaurada , sã como a outra.
\VS{6}Assim que os fariseus saíram, tiveram reunião com os herodianos contra ele, para combinarem sobre como o matariam.
\VS{7}E Jesus retirou-se com os seus discípulos para o mar; e seguiu-o uma grande multidão da Galileia, da Judeia,
\VS{8}de Jerusalém, da Idumeia, dalém do Jordão, e os das proximidades de Tiro e de Sidom; uma grande multidão, tendo ouvido quão grandes coisas fazia, vieram a ele.
\VS{9}E disse aos seus discípulos que um barquinho ficasse continuamente perto dele, por causa das multidões; para que não o apertassem.
\VS{10}Pois havia curado muitos, de maneira que todos quantos tinham algum mal lançavam-se sobre ele a fim de tocá-lo.
\VS{11}E os espíritos imundos, quando o viam, prostravam-se diante dele, e exclamavam: Tu és o Filho de Deus.
\VS{12}Mas
{\ADD{Jesus}} os repreendia muito, para que não manifestassem quem ele era.
\VS{13}Ele subiu ao monte, e chamou para si os que quis; então vieram a ele.
\VS{14}E constituiu doze para que estivessem com ele, para enviá-los a pregar,
\VS{15}para que tivessem poder de curar as enfermidades e expulsar os demônios.
\VS{16}Eram eles: Simão, a quem pôs por nome Pedro;
\VS{17}Tiago
{\ADD{filho}} de Zebedeu, e João, irmão de Tiago; e pôs-lhes por nome Boanerges, que significa “filhos do trovão”;
\VS{18}e André, Filipe, Bartolomeu, Mateus, Tomé; Tiago
{\ADD{filho}} de Alfeu; Tadeu; Simão o zelote;
\FTNT{*}{{\FR 3:18 }
Tradicionalmente “cananeu”, mas provavelmente o significado não tenha a ver com Caná ou Canaã
}
\VS{19}e Judas Iscariotes, o que o traiu.
\VS{20}Quando foram para uma casa, outra vez se ajuntou uma multidão, de maneira que nem sequer podiam comer pão.
\VS{21}Os seus
{\ADD{familiares}} , ao ouvirem isso, saíram para detê-lo, porque diziam: “Ele stá fora de si”.
\VS{22}E os escribas que haviam descido de Jerusalém diziam: Ele tem Belzebu, e é pelo chefe dos demônios que expulsa demônios.
\VS{23}Então
{\ADD{Jesus}} os chamou, e lhes disse por parábolas: Como pode Satanás expulsar Satanás?
\VS{24}Se algum reino estiver dividido contra si mesmo, tal reino não pode durar;
\VS{25}e se alguma casa estiver dividida contra si mesma, tal casa não pode durar firme.
\VS{26}E se Satanás se levantar contra si mesmo, e for dividido, não pode durar, mas tem fim.
\VS{27}Ninguém pode roubar os bens do valente, quando se entra na casa dele, se antes não amarrar ao valente; depois disso roubará a sua casa.
\VS{28}Em verdade vos digo que todos os pecados serão perdoados aos filhos dos homens, e todas as blasfêmias com que blasfemarem;
\VS{29}mas quem blasfemar contra o Espírito Santo ficará sem perdão para sempre; em vez disso, é culpado do juízo eterno.
\VS{30}Pois diziam: “Ele tem espírito imundo”.
\VS{31}Então chegaram os seus irmãos e a sua mãe ; e estando de fora, mandaram chamá-lo.
\VS{32}A multidão estava sentada ao redor dele. Então disseram-lhe: Eis que a tua mãe e os teus irmãos estão lá fora a te procurar.
\VS{33}Ele lhes respondeu: Quem é a minha mãe ou os meus irmãos?
\VS{34}E, olhando em redor aos que estavam sentados perto dele, disse: Eis aqui minha mãe e meus irmãos.
\VS{35}Pois quem fizer a vontade de Deus, esse é meu irmão, minha irmã, e mãe.

\par }\Chap{4}{\PP \VerseOne{1}{\ADD{Jesus}} começou outra vez a ensinar junto ao mar, e uma grande multidão se ajuntou a ele, de maneira que ele entrou num barco e ficou sentado no mar; e toda a multidão estava em terra junto ao mar.
\VS{2}E ensinava-lhes muitas coisas por parábolas; e dizia-lhes em seu ensinamento:
\VS{3}Ouvi: eis que o semeador saiu a semear;
\VS{4}E aconteceu que, enquanto semeava, uma
{\ADD{parte das sementes}} caiu junto ao caminho, e os pássaros do céu vieram, e a comeram.
\VS{5}E outra caiu em pedregulhos, onde não havia muita terra; e logo nasceu, porque não tinha terra profunda.
\VS{6}Mas, quando saiu o sol, queimou-se; e porque não tinha raiz, secou-se.
\VS{7}E outra caiu entre espinhos; e os espinhos cresceram e a sufocaram, e não deu fruto.
\VS{8}Mas outra caiu em boa terra, e deu fruto, que subiu, e cresceu; e um deu trinta, outro sessenta, e outro cem.
\VS{9}E disse-lhes : Quem tem ouvidos para ouvir, ouça.
\VS{10}E quando
{\ADD{Jesus}} esteve só, os que estavam junto dele, com os doze, perguntaram-lhe acerca da parábola .
\VS{11}E respondeu-lhes: A vós é concedido saber o mistério do Reino de Deus; mas aos que são de fora, todas estas coisas se fazem por meio de parábolas;
\VS{12}para que vendo, vejam, e não percebam; e ouvindo, ouçam, e não entendam; para não haver de se converterem, e lhes sejam perdoados os pecados .
{\RQ{Isaías 6:9-10
}}
\VS{13}E disse-lhes: Não sabeis o significado desta parábola? Como, pois, entendereis todas as parábolas?
\VS{14}O semeador semeia a palavra.
\VS{15}E estes são os de junto ao caminho: nos quais a palavra é semeada; mas depois de a ouvirem, Satanás logo vem, e tira a palavra que foi semeada nos seus corações .
\VS{16}E, semelhantemente, estes são os que se semeiam em pedregulhos: os que havendo ouvido a palavra, logo a recebem com alegria.
\VS{17}Mas não têm raiz em si mesmos; em vez disso, são temporários. Depois, quando se levanta a tribulação ou a perseguição por causa da palavra, logo tropeçam na fé.
\VS{18}E estes são os que se semeiam entre espinhos: os que ouvem a palavra;
\VS{19}mas as preocupações deste mundo, a sedução das riquezas, e as cobiças por outras coisas, entram, sufocam a palavra, e ela fica sem gerar fruto.
\VS{20}E estes são os que foram semeados em boa terra: os que ouvem a palavra, recebem-na, e dão fruto, um trinta, e outro sessenta, e outro cem.
\VS{21}E ele lhes disse: Por acaso a lâmpada vem a ser posta debaixo de uma caixa ou sob da cama? Não deve ela ser posta na luminária?
\VS{22}Pois não há nada encoberto que não haja de ser revelado; e nada se faz
{\ADD{para ficar}} encoberto, mas sim, para ser vir à luz.
\VS{23}Se alguém tem ouvidos para ouvir, ouça.
\VS{24}E disse-lhes: Prestai atenção ao que ouvis: com a medida que medirdes a vós mesmos se medirá, e será acrescentado a vós que ouvis .
\VS{25}Pois ao que tem, lhe será dado; e ao que não tem, até o que tem lhe será tirado.
\VS{26}E dizia: Assim é o Reino de Deus, como se um homem lançasse semente na terra;
\VS{27}e dormisse, e se levantasse, de noite e de dia, e a semente brotasse, e crescesse, sem que ele saiba como.
\VS{28}Pois a terra de si mesma frutifica, primeiro a erva, depois a espiga, depois o grão cheio na espiga.
\VS{29}E quando o fruto se mostra pronto, logo mete a foice, pois a colheita chegou.
\VS{30}E dizia: A que assemelharemos o Reino de Deus? Ou com que parábola o compararemos?
\VS{31}Com um grão da mostarda que, quando semeado na terra, é a menor de todas as sementes na terra.
\VS{32}Mas, depois de semeado, cresce, e se torna a maior de todas as hortaliças, e cria grandes ramos, de maneira que os pássaros do céu podem fazer ninhos sob a sua sombra.
\VS{33}E com muitas parábolas como essas
{\ADD{Jesus
}} lhes falava a palavra, conforme o que podiam ouvir.
\VS{34}E não lhes falava sem parábola; mas aos seus discípulos explicava tudo em particular.
\VS{35}Naquele dia, chegando o entardecer, disse-lhes: Passemos para o outro lado.
\VS{36}Então despediram a multidão, e o levaram consigo assim como estava no barco; mas havia também outros barquinhos com ele.
\VS{37}E levantou-se uma grande tempestade de vento; as ondas atingiam por cima do barco, de maneira que já se enchia.
\VS{38}{\ADD{Jesus}} estava na popa dormindo sobre uma almofada. Então despertaram-no, e disseram-lhe: Mestre, não te importas que pereçamos?
\VS{39}Então ele se levantou, repreendeu o vento, e disse ao mar: Cala-te, aquieta-te! E o vento se aquietou, e fez-se grande bonança.
\VS{40}E perguntou-lhes: Por que sois tão covardes? Como não tendes fé?
\VS{41}E ficaram muito atemorizados, e diziam uns aos outros: Quem é este, que até o vento e o mar lhe obedecem?

\par }\Chap{5}{\PP \VerseOne{1}E chegaram ao outro lado do mar, à terra dos Gadarenos.
\VS{2}Assim que
{\ADD{Jesus}} saiu do barco, veio das sepulturas ao seu encontro um homem com um espírito imundo,
\VS{3}que morava nas sepulturas, e nem mesmo com correntes conseguiam prendê-lo;
\VS{4}pois muitas vezes fora preso com grilhões e correntes; mas as correntes eram por ele feitas em pedaços, os grilhões eram esmigalhados, e ninguém o conseguia controlar.
\VS{5}E sempre dia e noite andava gritando pelos montes e pelas sepulturas, e ferindo-se com pedras.
\VS{6}Quando ele viu Jesus de longe, correu e prostrou-se diante dele.
\VS{7}E gritou em alta voz: Que tenho eu contigo Jesus, Filho do Deus Altíssimo? Imploro-te por Deus que não me atormentes.
\VS{8}(Pois
{\ADD{Jesus}} havia lhe dito: “Sai deste homem, espírito imundo”.)
\VS{9}Então perguntou-lhe: Qual é o teu nome? E respondeu: Legião é o meu nome, porque somos muitos.
\VS{10}E rogava-lhe muito que não os expulsasse daquela terra.
\VS{11}Havia ali perto dos montes uma grande manada de porcos pastando.
\VS{12}E todos
{\ADD{aqueles
}} demônios rogaram-lhe, dizendo: Manda-nos para aqueles porcos, para que entremos neles.
\VS{13}Imediatamente Jesus lhes permitiu. Então aqueles espíritos imundos saíram para entrar nos porcos; e a manada lançou-se abaixo no mar; (eram quase dois mil) e afogaram-se no mar.
\VS{14}Os que apascentavam os porcos fugiram, e avisaram na cidade e nos campos; e
{\ADD{as pessoas}} foram ver o que havia acontecido.
\VS{15}Então aproximaram-se de Jesus, e viram o endemoninhado sentado, vestido, e em sã consciência o que tivera a legião; e ficaram apavorados.
\VS{16}E os que haviam visto contaram-lhes o que acontecera ao endemoninhado, e sobre os porcos.
\VS{17}Então começaram a rogar-lhe que saísse do território deles.
\VS{18}Quando
{\ADD{Jesus}} entrava no barco, o que fora endemoninhado rogou-lhe que estivesse com ele.
\VS{19}Jesus se recusou, porém lhe disse: Vai para a tua casa, aos teus, e anuncia-lhes quão grandes coisas o Senhor fez contigo, e
{\ADD{como}} teve misericórdia de ti.
\VS{20}Então ele foi embora, e começou a anunciar em Decápolis quão grandes coisas Jesus havia feito com ele; e todos se admiravam.
\VS{21}Depois de Jesus passar outra vez num barco para o outro lado, uma grande multidão se ajuntou a ele; e ele ficou junto ao mar.
\VS{22}E eis que veio um dos líderes de sinagoga, por nome Jairo; e quando o viu, prostrou-se aos seus pés.
\VS{23}E implorava-lhe muito, dizendo: Minha filhinha está a ponto de morrer.
{\ADD{Rogo-te}} que venhas pôr as mãos sobre ela, para que seja curada, e viva.
\VS{24}{\ADD{Jesus}} foi com ele. Uma grande multidão o seguia, e o apertavam.
\VS{25}E havia uma certa mulher, que tinha um fluxo de sangue havia doze anos,
\VS{26}que tinha sofrido muito por meio de muitos médicos, e gastado tudo quanto possuía, e nada havia lhe dado bom resultado; ao invés disso, piorava.
\VS{27}Quando ela ouviu falar de Jesus, veio entre a multidão por detrás, e tocou a roupa dele.
\VS{28}Pois dizia: Se tão somente tocar as suas roupas, serei curada.
\VS{29}E imediatamente a fonte do seu sangue parou se secou; e sentiu no corpo que já havia sido curada daquele flagelo.
\VS{30}Jesus logo notou em si o poder que dele havia saído. Então virou-se na multidão, e perguntou: Quem tocou as minhas roupas?
\VS{31}E seus discípulos lhe disseram: Eis que a multidão te aperta, e perguntas: Quem me tocou?
\VS{32}E ele olhava em redor, para ver quem havia lhe feito isso.
\VS{33}Então a mulher temendo, e tremendo, sabendo o que havia sido feito em si, veio, prostrou-se diante dele, e disse-lhe toda a verdade.
\VS{34}E ele lhe disse: Filha, a tua fé te salvou. Vai em paz, e estejas curada deste teu flagelo.
\VS{35}Estando ele ainda falando, alguns vieram
{\ADD{da casa}} do líder de sinagoga, e disseram: A tua filha já morreu; por que ainda estás incomodando o Mestre?
\VS{36}Mas Jesus, assim que ouviu essa palavra que havia sido falada, disse ao líder de sinagoga: Não temas; crê somente.
\VS{37}E não permitiu que ninguém o seguisse, a não ser Pedro, Tiago, e João irmão de Tiago.
\VS{38}Ele chegou à casa do líder de sinagoga, e viu o alvoroço, os que choravam muito e pranteavam.
\VS{39}E ao entrar, disse-lhes: Por que fazeis alvoroço e chorais? A menina não morreu, mas está dormindo.
\VS{40}E riram dele. Porém ele, depois de pôr todos fora, tomou consigo o pai e a mãe da menina, e os que estavam com ele. Em seguida, entrou onde a menina estava deitada.
\VS{41}Ele pegou a mão da menina, e lhe disse: “Talita cumi”, (que significa: “Menina, eu te digo, levanta-te”).
\VS{42}E logo a menina se levantou e andou, pois já tinha doze anos de idade. E ficaram grandemente espantados.
\VS{43}E mandou-lhes muito que ninguém o soubesse; e mandou que dessem a ela de comer.

\par }\Chap{6}{\PP \VerseOne{1}{\ADD{Jesus}} partiu-se dali, veio à sua terra, e seus discípulos o seguiram.
\VS{2}E chegando o sábado, começou a ensinar na sinagoga; e muitos, quando o ouviram, espantavam-se, dizendo: De onde lhe
{\ADD{vem}} estas coisas? E que sabedoria é esta que lhe foi dada? E tais maravilhas feitas por suas mãos?
\VS{3}Não é este o carpinteiro, filho de Maria, e irmão de Tiago, de José, de Judas, e de Simão? E não estão aqui as suas irmãs conosco? E ofenderam-se nele.
\VS{4}E Jesus lhes dizia: Todo profeta tem honra,
\FTNT{*}{{\FR 6:4 }
Lit. “Não há profeta sem honra”
} menos em sua terra, entre os parentes, e em sua própria casa.
\VS{5}Ele não pôde ali fazer milagre algum, a não ser somente para uns poucos enfermos, sobre os quais pôs as mãos e os curou.
\VS{6}E ficou admirado da incredulidade deles. Ele percorreu as aldeias do redor, ensinando.
\VS{7}E chamou a si os doze, e começou a enviar de dois em dois; e deu-lhes poder sobre os espíritos imundos.
\VS{8}E mandou-lhes que não tomassem nada para o caminho, a não ser somente um bordão; nem bolsa, nem pão, nem dinheiro no cinto;
\VS{9}mas que calçassem sandálias, e não se vestissem de duas túnicas.
\VS{10}E dizia-lhes: Onde quer que entrardes em alguma casa, ficai ali até que dali saiais.
\VS{11}E todos os que não vos receberem, nem vos ouvirem, quando sairdes dali, sacudi o pó que estiver debaixo de vossos pés, em testemunho contra eles Em verdade vos digo, que mais tolerável será a
{\ADD{os de}} Sodoma ou Gomorra no dia do juízo, do que a
{\ADD{os d}} aquela cidade.
\VS{12}Eles, então, se foram, e pregaram que
{\ADD{as pessoas}} se arrependessem.
\VS{13}Eles expulsaram muitos demônios, e a muitos enfermos ungiram com azeite, e os curaram.
\VS{14}O rei Herodes ouviu falar disso (porque o nome de
{\ADD{Jesus}} já era notório). E dizia: João Batista ressuscitou dos mortos, e por isso estas maravilhas operam nele.
\VS{15}Outros diziam: É Elias; e outros diziam: É profeta, ou como algum dos profetas.
\VS{16}Quando, porém, Herodes ouviu falar disso, falou: Ele é João, de quem cortei a cabeça. Ele ressuscitou dos mortos.
\VS{17}Pois o próprio Herodes havia mandado prender João, e acorrentá-lo na prisão, por causa de Herodias, mulher do seu irmão Filipe, porque havia se casado com ela.
\VS{18}Pois João dizia a Herodes: Não te é lícito possuir a mulher do teu irmão.
\VS{19}Assim Herodias o odiava, e queria matá-lo, mas não podia,
\VS{20}pois Herodes temia João, sabendo que era um homem justo e santo, e o estimava. E quando o ouvia, fazia muitas coisas, o ouvia de boa vontade.
\VS{21}Mas veio um dia oportuno, em que Herodes, no dia do seu aniversário, dava uma ceia aos grandes de sua corte, aos comandantes militares, e aos principais da Galileia.
\VS{22}Então a filha dessa Herodias entrou dançando, e agradou a Herodes e aos que estavam sentados com ele. O rei disse à garota: Pede-me quanto quiseres, que eu darei a ti.
\VS{23}E jurou a ela: Tudo o que me pedirdes te darei, até a metade do meu reino.
\VS{24}Então ela saiu, e perguntou à sua mãe: Que pedirei? E ela respondeu: A cabeça de João Batista.
\VS{25}E entrando ela logo apressadamente ao rei, pediu, dizendo: Quero que imediatamente me dês num prato a cabeça de João Batista.
\VS{26}E o rei entristeceu-se muito;
{\ADD{mas}} , por causa dos juramentos, e dos que estavam juntamente à mesa, não quis recusar a ela.
\VS{27}Então logo o rei enviou o executor com a ordem de trazer ali sua cabeça. Ele, foi, e o decapitou na prisão.
\VS{28}Em seguida, trouxe a sua cabeça num prato, e o deu à garota; e a garota a deu à sua mãe.
\VS{29}Quando os discípulos dele ouviram isso, vieram, pegaram o seu cadáver, e o puseram num sepulcro.
\VS{30}Os apóstolos juntaram-se
{\ADD{de volta}} a Jesus, e contaram-lhe tudo, tanto o que haviam feito, como o que haviam ensinado.
\VS{31}E ele lhes disse: Vinde vós à parte a um lugar deserto, e descansai um pouco; pois havia muitos que iam e vinham, e não tinham tempo para comer.
\VS{32}E foram-se num barco a um lugar deserto à parte.
\VS{33}Mas as multidões os viram ir, e muitos o reconheceram. Então correram para lá a pé de todas as cidades, chegaram antes deles, e vieram para perto dele.
\VS{34}Quando Jesus saiu
{\ADD{do barco}} , viu uma grande multidão, e teve compaixão deles porque eram como ovelhas que não têm pastor. Assim, começou a lhes ensinar muitas coisas.
\VS{35}E quando já era tarde, os seus discípulos vieram a ele, e disseram: O lugar é deserto, e a hora já é tarde.
\VS{36}Despede-os, para eles irem aos campos e aldeias circunvizinhos, e comprarem pão para si; pois não têm o que comer.
\VS{37}Mas ele respondeu: Dai-lhes vós mesmos de comer. E eles lhe responderam: Iremos, e compraremos duzentos denários de pão, para lhes darmos de comer?
\VS{38}E ele lhes disse: Quantos pães tendes? Ide ver. Quando souberam, disseram: Cinco, e dois peixes.
\VS{39}E mandou-lhes que fizessem sentar a todos em grupos sobre a grama verde.
\VS{40}E sentaram-se repartidos de cem em cem, e de cinquenta em cinquenta.
\VS{41}Ele tomou os cinco pães e os dois peixes, levantou os olhos ao céu, abençoou, e partiu os pães, e os deu aos seus discípulos, para que os pusessem diante deles. E os dois peixes repartiu com todos.
\VS{42}Todos comeram e se saciaram.
\VS{43}E dos pedaços de pão e dos peixes levantaram doze cestos cheios.
\VS{44}Os que comeram os pães eram quase cinco mil homens.
\VS{45}Logo depois, ordenou seus discípulos a subirem no barco, e ir adiante para o outro lado, em Betsaida, enquanto ele despedia a multidão.
\VS{46}E, depois de os despedir, foi ao monte para orar.
\VS{47}Ao anoitecer, o barco estava no meio do mar, e
{\ADD{Jesus}} sozinho em terra.
\VS{48}E viu que se cansavam muito remando, porque o vento lhes era contrário. Então, perto da quarta vigília da noite, veio a eles andando sobre o mar, e queria passar por eles.
\VS{49}Mas quando eles o viram andando sobre o mar, pensaram que era uma fantasma, e gritaram,
\VS{50}pois todos o viam, e ficaram perturbados. Então logo falou com eles, dizendo: Tende coragem! Sou eu, não tenhais medo.
\VS{51}E subiu a eles no barco, e o vento se aquietou. Eles ficaram muito espantados e maravilhados entre si,
\VS{52}pois não haviam entendido o que tinha acontecido com os pães, porque o coração deles estava endurecido.
\VS{53}Eles terminaram de atravessar o mar, chegaram à terra de Genesaré, e ali aportaram.
\VS{54}Quando eles saíram do barco, logo
{\ADD{as pessoas}} o reconheceram.
\VS{55}Então gente de toda a região em redor veio correndo, e começaram a trazer em camas os doentes, aonde quer que ouviam que ele estava.
\VS{56}E aonde quer que ele entrava, em povoados, cidades, ou aldeias, colocavam os enfermos nas praças, e rogavam-lhe que ao menos tocassem a borda de sua roupa; e todos os que o tocavam ficavam sarados.

\par }\Chap{7}{\PP \VerseOne{1}Os fariseus, e alguns dos escribas, que tinham vindo de Jerusalém, reuniram-se com ele.
\VS{2}E, quando viram que alguns dos discípulos dele comiam pão com mãos impuras, isto é, sem lavar, repreendiam-lhes.
\VS{3}(Pois os fariseus, e todos os judeus, mantendo a tradição dos antigos, se não lavarem bastante as mãos, não comem.
\VS{4}E, quando voltam da rua, se não se lavarem, não comem; e há muitas outras coisas que se encarregam de guardar,
{\ADD{como}} lavar os copos, as vasilhas, os utensílios de metal, e os leitos).
\VS{5}Depois os fariseus e os escribas lhe perguntaram: Por que os teus discípulos não andam conforme a tradição dos antigos, em vez de comerem pão com as mãos sem lavar?
\VS{6}E ele lhes respondeu: Bem profetizou Isaías acerca de vós, hipócritas! Como está escrito: Este povo me honra com os lábios, mas o seu coração está longe de mim.
\VS{7}Eles, porém, me honram em vão, ensinando como doutrinas mandamentos humanos.
{\RQ{Isaías 29:13
}}
\VS{8}Pois vós deixais o mandamento de Deus, e mantendes a tradição humana,
{\ADD{como}} lavar as vasilhas e os copos; e fazeis muitas outras coisas semelhantes a estas.
\VS{9}E dizia-lhes: Vós dispensais o mandamento de Deus para guardardes a vossa tradição;
\VS{10}porque Moisés disse: Honra o teu pai e a tua mãe. E quem maldisser ao pai ou à mãe terá de morrer.
\VS{11}Mas vós dizeis: Se o homem disser ao pai ou à mãe: Tudo o que te puder aproveitar de mim é corbã (isto é, oferta),
\VS{12}então não lhe deixais mais nada fazer por seu pai ou por sua mãe.
\VS{13}Assim invalidaias a palavra de Deus por vossa tradição, que ordenastes; e
{\ADD{fazeis}} muitas coisas semelhantes a estas.
\VS{14}E chamando para si toda a multidão, disse-lhes: Ouvi-me todos, e entendei:
\VS{15}Nada há fora do ser humano que nele entre que o possa contaminar; mas o que dele sai, isso é o que contamina o ser humano.
\VS{16}Se alguém tem ouvidos para ouvir, ouça.
\VS{17}Quando
{\ADD{Jesus}} deixou a multidão e entrou em casa, seus discípulos lhe perguntaram sobre a parábola.
\VS{18}E ele lhes disse: Também vós estais assim sem entendimento? Não entendeis que tudo o que de fora entra no ser humano não o pode contaminar?
\VS{19}Pois não entra no seu coração, mas, sim, no ventre, e sai para a privada. (Assim, ele declarou como limpas todas as comidas).
\VS{20}E dizia: O que sai do ser humano, isso contamina o ser humano.
\VS{21}Pois é de dentro do coração humano que vêm os maus pensamentos, os adultérios, os pecados sexuais, os homicídios,
\VS{22}os roubos, as ganâncias, as maldades, o engano, a depravação, o olho malicioso, a blasfêmia, a soberba, a insensatez.
\VS{23}Todos estes males procedem de dentro, e contaminam o ser humano.
\VS{24}{\ADD{Jesus}} levantou-se dali e foi para a região de Tiro e de Sidom. Ele entrou numa casa, e não queria que ninguém soubesse disso, mas não pôde se esconder.
\VS{25}Pois uma mulher, cuja filhinha tinha um espírito imundo, assim que ouviu falar dele, veio, e prostrou-se a seus pés.
\VS{26}Esta mulher era grega, de nacionalidade sirofenícia; e rogava-lhe que expulsasse o demônio de sua filha.
\VS{27}Mas Jesus lhe disse: Deixa primeiro que os filhos se fartem; porque não é bom tomar o pão dos filhos, e lançá-lo aos cachorrinhos.
\VS{28}Porém ela lhe respondeu: Sim Senhor; mas também os cachorrinhos comem debaixo da mesa, das migalhas que os filhos deixam.
\VS{29}Então ele lhe disse: Por esta palavra, vai, o demônio já saiu da tua filha.
\VS{30}Quando ela chegou à sua casa, encontrou que o demônio já havia saído, e a filha estava deitada sobre a cama.
\VS{31}Então
{\ADD{Jesus}} voltou a sair da região de Tiro e de Sidom, e veio para o mar da Galilea, por meio da região de Decápolis.
\VS{32}E trouxeram-lhe um surdo que dificilmente falava, e rogaram-lhe que pusesse a mão sobre ele.
\VS{33}E tomando-o em separado da multidão, pôs os seus dedos nos ouvidos dele, cuspiu, e tocou-lhe a língua.
\VS{34}Depois, levantando os olhos ao céu, suspirou, e disse: Efatá, (isto é, abre-te).
\VS{35}Imediatamente os ouvidos dele se abriram, e o que prendia sua língua se soltou, e passou a falar bem.
\VS{36}{\ADD{Jesus}} lhes mandou que a ninguém dissessem; porém, quanto mais lhes mandava, mais divulgavam.
\VS{37}E ficavam muito admirados, dizendo: Ele faz tudo bem! Aos surdos faz ouvir, e aos mudos falar.

\par }\Chap{8}{\PP \VerseOne{1}Naqueles dias, quando havia uma multidão muito grande, e não tinham o que comer, Jesus chamou os seus discípulos a si, e disse-lhes:
\VS{2}Tenho compaixão da multidão, porque já há três dias que estão comigo, e não têm o que comer.
\VS{3}E se eu os deixar ir sem comer para suas casas desmaiarão no caminho; porque alguns deles vieram de longe.
\VS{4}Os seus discípulos lhe responderam: De onde poderá alguém saciar estes de pão aqui no deserto?
\VS{5}{\ADD{Jesus}} lhes perguntou: Quantos pães tendes? E eles disseram: Sete.
\VS{6}Então mandou à multidão que se sentassem no chão. Em seguida, tomou os sete pães, e tendo dado graças, partiu-os, e os deu a seus discípulos, para que os pusessem diante deles. E eles os puseram diante da multidão.
\VS{7}E tinham uns poucos peixinhos, e havendo dado graças, disse que também os pusessem diante deles.
\VS{8}Eles comeram, e se saciaram; e levantaram, do que sobrou dos pedaços, sete cestos.
\VS{9}Os que comeram eram quase quatro mil. Depois os despediu.
\VS{10}E logo entrou no barco com os seus discípulos, e veio para a região de Dalmanuta.
\VS{11}Os fariseus vieram, e começaram a disputar com ele, pedindo-lhe sinal do céu, para o testar.
\VS{12}E ele, suspirando profundamente em seu espírito, disse: Por que esta geração pede um sinal? Em verdade vos digo, que não se dará sinal a esta geração.
\VS{13}Então os deixou, voltou a entrar no barco , e foi para a outro lado
{\ADD{do mar}} .
\VS{14}E
{\ADD{os seus discípulos}} haviam se esquecido de tomar pão, e nada tinham, a não ser um pão com eles no barco.
\VS{15}E
{\ADD{Jesus}} lhes deu a seguinte ordem: Prestai atenção: tende cuidado com o fermento dos fariseus e o fermento de Herodes.
\VS{16}E indagavam-se com os outros , dizendo: É porque não temos pão.
\VS{17}Jesus soube e lhes disse: Por que indagais que não tendes pão? Não percebeis ainda, nem entendeis? Ainda tendes o vosso coração endurecido?
\VS{18}Tendes olhos, e não vedes? Tendes ouvidos, e não ouvis?
\VS{19}E não vos lembrais de, quando parti os cinco pães entre os cinco mil, quantos cestos cheios de pedaços levantastes?Responderam-lhe: Doze.
\VS{20}E quando
{\ADD{parti}} os sete entre os quatro mil, quantas cestas cheias de pedaços levantastes? Eles disseram: Sete.
\VS{21}Ele lhes perguntou: Como não entendeis?
\VS{22}Então veio a Betsaida. E trouxeram-lhe um cego, e rogaram-lhe que o tocasse.
\VS{23}Ele tomou o cego pela mão e o tirou para fora da aldeia. Depois cuspiu nos olhos dele e, pondo as mãos encima dele, perguntou-lhe se via alguma coisa.
\VS{24}Ele levantou os olhos e disse: Vejo as pessoas; pois vejo como árvores que andam.
\VS{25}Então
{\ADD{Jesus}} pôs de novo as mãos sobre os seus olhos, e o fez olhar para cima . Assim ele ficou restabelecido, e passou a ver todos claramente.
\VS{26}Então o mandou para sua casa, dizendo: Não entres na aldeia, nem contes a ninguém da aldeia.
\VS{27}Jesus saiu com os seus discípulos para as aldeias de Cesareia de Filipe; e no caminho perguntou a seus discípulos: Quem as pessoas dizem que eu sou?
\VS{28}Eles responderam: João Batista, e outros, Elias; e outros, algum dos profetas.
\VS{29}E ele lhes perguntou: E vós, quem dizeis que eu sou? Pedro lhe respondeu: Tu és o Cristo.
\VS{30}E lhes ordenou que a ninguém dissessem aquilo dele.
\VS{31}E começou a lhes ensinar que era necessário que o Filho do homem sofresse muito, e fosse rejeitado pelos anciãos e pelos chefes dos sacerdotes e escribas, e que fosse morto, e depois de três dias ressuscitasse.
\VS{32}Ele dizia essa palavra abertamente. Então Pedro o tomou à parte, e começou a repreendê-lo.
\VS{33}Mas
{\ADD{Jesus}} virou-se e, olhando para os seus discípulos, repreendeu Pedro, dizendo: Sai de diante de mim, satanás! Pois tu não compreendes as coisas de Deus, mas sim as humanas.
\VS{34}Então chamou a si a multidão com os seus discípulos, e disse-lhes: Quem quiser vir após mim, negue a si mesmo, tome a sua cruz, e siga-me;
\VS{35}pois quem quiser salvar a sua vida, a perderá; mas quem perder a sua vida por causa de mim e do Evangelho, esse a salvará.
\VS{36}Pois que proveito teria alguém, se ganhasse o mundo todo, e perdesse a sua alma?
\VS{37}Ou que daria alguém em resgate de sua alma?
\VS{38}Porque todo aquele que se envergonhar de mim e de minhas palavras nesta geração adúltera e pecadora, também o Filho do homem, quando vier na glória de seu Pai com os santos anjos, envergonhar-se-á dele.

\par }\Chap{9}{\PP \VerseOne{1}E disse-lhes também: Em verdade vos digo, que há alguns dos que aqui estão, que não experimentarão a morte, até que vejam o reino de Deus vindo com poder.
\VS{2}Seis dias depois, Jesus tomou consigo Pedro, Tiago, e João, e os levou à parte, sozinhos, para um alto monte ; e transfigurou-se diante deles.
\VS{3}E suas roupas ficaram resplandescentes, muito brancas como a neve , como nenhum lavadeiro na terra seria capaz de branquear.
\VS{4}E apareceu-lhes Elias com Moisés, e falavam com Jesus.
\VS{5}Então Pedro disse a Jesus: Mestre, é bom para nós estarmos aqui; façamos três tendas: uma para ti, uma para Moisés, e uma para Elias.
\VS{6}Pois ele não sabia o que dizia, pois estavam assombrados.
\VS{7}Então desceu uma nuvem, que os cobriu com a sua sombra, e veio uma voz da nuvem, que dizia: Este é meu Filho amado; a ele ouvi.
\VS{8}De repente, quando olharam em redor, não viram mais ninguém, a não ser só Jesus com eles.
\VS{9}Enquanto desciam do monte,
{\ADD{Jesus}} lhes mandou que a ninguém contassem o que haviam visto, até que o Filho do homem ressuscitasse dos mortos.
\VS{10}E eles guardaram o caso entre si, perguntando uns aos outros o que seria aquilo de “ressuscitar dos mortos”.
\VS{11}Então lhe perguntaram: Por que os escribas dizem que Elias tem que vir primeiro?
\VS{12}E ele lhes respondeu: De fato Elias vem primeiro, e restaura todas as coisas. Então, como está escrito sobre o Filho do homem tem que sofrer muito, e ser desprezado?
\VS{13}Porém eu vos digo que Elias já veio, e fizeram-lhe tudo o que quiseram, como está escrito sobre ele.
\VS{14}E quando veio aos discípulos, ele viu uma grande multidão ao redor deles; e uns escribas estavam discutindo com eles.
\VS{15}Logo que toda a multidão o viu, ficou admirada. Então correram a ele, e o cumprimentaram.
\VS{16}{\ADD{Jesus}} perguntou aos escribas: O que estais discutindo com eles?
\VS{17}E um da multidão respondeu: Mestre, trouxe a ti o meu filho, que tem um espírito mudo.
\VS{18}E onde quer que o toma, faz-lhe ter convulsões, solta espuma, range os dentes, e vai ficando rígido. Eu disse aos teus discípulos que o expulsassem, mas não conseguiram.
\VS{19}{\ADD{Jesus}} lhe respondeu: Ó geração incrédula! Até quando estarei ainda convosco? Até quando vos suportarei? Trazei-o a mim.
\VS{20}Então trouxeram-no a ele. E quando o viu, logo o espírito o fez ter uma convulsão e, caindo em terra, rolava, e espumava.
\VS{21}E perguntou ao seu pai: Quanto tempo há que isto lhe sobreveio? E ele lhe disse: Desde a infância.
\VS{22}E muitas vezes o lançou também no fogo e na água para o destruir. Mas, se podes algo, tem compaixão de nós, e ajuda-nos.
\VS{23}E Jesus lhe disse: Se podes crer, tudo é possível ao que crê.
\VS{24}Logo o pai do menino, clamando, com lágrimas, disse: Creio, Senhor! Ajuda-me na minha incredulidade.
\VS{25}Quando Jesus viu que a multidão corria e se ajuntava, repreendeu o espírito imundo, dizendo-lhe: Espírito mudo e surdo, eu te ordeno, sai dele, e nunca mais entres nele!
\VS{26}Então
{\ADD{o espírito}} , gritando, e fazendo-o ter muita convulsão, saiu. E
{\ADD{o menino}} ficou como morto, de maneira que muitos diziam que estava morto.
\VS{27}Então Jesus o tomou pela mão, e o ergueu; e ele se levantou.
\VS{28}E quando entrou em casa, seus discípulos lhe perguntaram à parte: Por que nós não conseguimos o expulsar?
\VS{29}E lhes respondeu: Este tipo não pode sair por coisa alguma, a não ser com oração e jejum.
\VS{30}Depois partiram dali, e caminharam pela Galileia. Mas
{\ADD{Jesus}} não queria que ninguém soubesse,
\VS{31}Porque ensinava a seus discípulos, e lhes dizia: O Filho do homem será entregue em mãos de homens, e o matarão; e
{\ADD{estando}} ele morto, ressuscitará ao terceiro dia.
\VS{32}Mas eles não entendiam esta palavra, e temiam lhe perguntar.
\VS{33}E veio a Cafarnaum, e entrando em casa, perguntou-lhes: Que indagais entre vós pelo caminho?
\VS{34}Mas eles se calaram; porque eles haviam discutido uns com os outros pelo caminho, qual
{\ADD{deles seria}} o maior.
\VS{35}E sentando-se ele, chamou aos doze, e disse-lhes: Se alguém quiser ser o primeiro, seja o últimos de todos, e servo de todos.
\VS{36}E tomando um menino, ele o pôs no meio deles, e tomando-o entre seus braços, disse-lhes:
\VS{37}Qualquer que em meu nome receber a um dos tais meninos, recebe a mim; e qualquer que me receber, não é a mim que recebe, mas, sim, àquele que me enviou.
\VS{38}E respondeu-lhe João, dizendo: Mestre, vimos um que em teu nome expulsava os demônios, e ele não nos segue. Então o proibimos, porque não nos segue.
\VS{39}Porém Jesus disse: Não o proibais; porque ninguém há que faça milagre em meu nome, e logo possa falar mal de mim.
\VS{40}Pois quem não é contra nós, é por nós.
\VS{41}Porque qualquer que vos der um copo d'água para beber em meu nome, porque sois de Cristo, em verdade vos digo, que não perderá sua recompensa.
\VS{42}E qualquer que conduzir ao pecado
\FTNT{*}{{\FR 9:42 }
conduzir ao pecado
Lit. fizer tropeçar
} um destes pequenos que creem em mim, melhor lhe fora que lhe pusesse ao pescoço uma grande pedra de moinho, e que fosse lançado no mar.
\VS{43}E se a tua mão te faz pecar,
\FTNT{†}{{\FR 9:43 }
Lit. “tropeçar”. tradicionalmente “escandalizar”. Também nos vv. 45 e 47
} corta-a; melhor te é entrar na vida mutilado, do que, tendo duas mãos, ir ao inferno, ao fogo que nunca se apaga.
\VS{44}Onde seu verme não morre, e o fogo nunca se apaga.
\VS{45}E se teu pé te faz pecar, corta-o; melhor te é entrar na vida manco, do que, tendo dois pés, ser lançado no inferno, no fogo que nunca se apaga.
\VS{46}Onde seu verme não morre, e o fogo nunca se apaga.
\VS{47}E se teu olho te faz pecar, lança-o fora; melhor te é entrar no Reino de Deus com um olho, do que, tendo dois olhos, ser lançado no fogo do inferno,
\VS{48}onde seu verme não morre, e o fogo nunca se apaga.
\VS{49}Porque cada um será salgado com fogo, e cada sacrifício será salgado com sal.
\VS{50}O sal é bom; mas se o sal se tornar insípido, com que o temperareis? Tende sal em vós mesmos, e paz uns com os outros.

\par }\Chap{10}{\PP \VerseOne{1}E levantando-se dali,
{\ADD{Jesus}} se foi aos limites da Judeia, por dalém do Jordão; e as multidões voltaram a se juntar a ele, e voltou a lhes ensinar, como tinha costume.
\VS{2}E vindo a ele os fariseus, perguntaram-lhe se era lícito ao homem deixar a
{\ADD{sua}} mulher, testando-o.
\VS{3}Mas ele lhes respondeu: O que Moisés vos mandou?
\VS{4}E eles disseram: Moisés permitiu escrever carta de divórcio, e deixá-la.
\VS{5}E Jesus lhes respondeu: Foi pela dureza dos vossos corações que ele vos escreveu este mandamento.
\VS{6}Porém, desde o princípio da criação, macho e fêmea Deus os fez.
{\RQ{Gênesis 1:27
}}
\VS{7}Por isso, deixará o homem a seu pai e a
{\ADD{sua}} mãe, e se unirá à sua mulher.
\VS{8}E os dois serão uma
{\ADD{só}} carne; assim, já não são dois, mas sim uma
{\ADD{só}} carne.
\VS{9}Portanto, o que Deus juntou, não separe o homem.
\VS{10}E em casa seus discípulos voltaram a lhe perguntar sobre isto mesmo.
\VS{11}E disse-lhes: Qualquer um que deixar a sua mulher, e se casar com outra, adultera contra ela.
\VS{12}E se a mulher deixar o seu marido, e se casar com outro, adultera.
\VS{13}E lhe traziam crianças para que ele as tocasse, mas os discípulos repreendiam aos que
{\ADD{as}} traziam.
\VS{14}Porém Jesus, vendo, indignou-se, e lhe disse: Deixai vir as crianças a mim, e não as impeçais; porque delas é o Reino de Deus.
\VS{15}Em verdade vos digo, que qualquer um que não receber o Reino de Deus como criança, em maneira nenhuma nele entrará.
\VS{16}Então, tomando-as entre seus braços, e pondo as mãos sobre elas, ele as abençoou.
\VS{17}E quando ele saiu para caminhar, um
{\ADD{homem}} correu até ele; e pondo-se de joelhos diante dele, perguntou-lhe: Bom Mestre, que farei para herdar a vida eterna?
\VS{18}E Jesus lhe disse: Por que me chamas bom? Ninguém é bom, a não ser um: Deus.
\VS{19}Sabes os mandamentos: Não adulterarás; não matarás; não furtarás; não darás falso testemunho; não serás enganador; honra o teu pai e a
{\ADD{tua}} mãe.
\VS{20}Porém ele lhe respondeu: Mestre, tudo isto guardei desde a minha juventude.
\VS{21}E olhando Jesus para ele, amou-o, e disse-lhe: Uma coisa te falta: vai, vende tudo quanto tens, e dá aos pobres; e terás um tesouro no céu; e vem, segue-me, toma
{\ADD{tua}} cruz.
\VS{22}Mas ele, pesaroso desta palavra, foi-se triste; porque tinha muitas propriedades.
\VS{23}Então Jesus olhando ao redor, disse a seus discípulos: Como é difícil aos que têm riquezas entrar no Reino de Deus!
\VS{24}E os discípulos se espantaram destas suas palavras; mas Jesus, voltando a responder, disse-lhes: Filhos, como é difícil aos que confiam em riquezas entrar no Reino de Deus!
\VS{25}Mais fácil é passar um camelo pelo olho de uma agulha, do que entrar um rico no Reino de Deus.
\VS{26}E eles se espantavam ainda mais, dizendo uns para os outros: Quem, pois, poderá se salvar?
\VS{27}Porém, olhando Jesus para eles, disse: Para os seres humanos, é impossível; mas para Deus, não; porque para Deus tudo é possível.
\VS{28}E Pedro começou a lhe dizer: Eis que nós deixamos tudo, e te seguimos.
\VS{29}Jesus respondeu: Em verdade vos digo, que não há ninguém que tenha deixado casa, ou irmãos, ou irmãs, ou pai, ou mãe, ou esposa, ou filhos, ou campos, por amor de mim e do Evangelho,
\VS{30}que não receba cem vezes tanto, agora neste tempo, casas, e irmãos, e irmãs, e mães, e filhos, e campos, com perseguições; e a vida eterna, no tempo que virá.
\VS{31}Porém muitos primeiros serão últimos, e últimos, primeiros.
\VS{32}E iam pelo caminho, subindo para Jerusalém; e Jesus ia diante deles, e espantavam-se, e seguiam-no atemorizados. E voltando a tomar consigo aos doze, começou-lhes a dizer as coisas que lhe viriam a acontecer:
\VS{33}Eis que subimos a Jerusalém, e o Filho do homem será entregue aos chefes dos sacerdotes, e aos escribas; e o condenarão à morte, e o entregarão aos gentios.
\VS{34}E escarnecerão dele, e o açoitarão, e cuspirão nele, e o matarão; e ao terceiro dia ressuscitará.
\VS{35}E vieram a ele Tiago e João, filhos de Zebedeu, dizendo: Mestre, queríamos que nos fizesses o que pedirmos.
\VS{36}E ele lhes disse: O que quereis que eu vos faça?
\VS{37}E eles lhe disseram: Concede-nos que em tua glória nos sentemos, um à tua direita, e outro à tua esquerda?
\VS{38}Mas Jesus lhes disse: Não sabeis o que pedis; podeis vós beber o copo que eu bebo, e ser batizados com o batismo com que eu sou batizado?
\VS{39}E eles lhe disseram: Podemos. Porém Jesus lhes disse: Em verdade, o copo que eu bebo, bebereis; e com o batismo com que eu sou batizado, sereis batizados.
\VS{40}Mas sentar-se à minha direita, e à minha esquerda, não é meu concedê-lo, mas sim, para aqueles a quem está preparado.
\VS{41}E quando os dez ouviram isto, começaram a se irritar contra Tiago e João.
\VS{42}Mas Jesus, chamando-os a si, disse-lhes: Já sabeis que os que são vistos como governadores dos gentios, agem como senhores deles; e os grandes usam de autoridade sobre eles.
\VS{43}Mas entre vós não será assim; antes qualquer que entre vós quiser ser grande, será vosso servidor.
\VS{44}E qualquer que de vós quiser ser o primeiro, será servo de todos.
\VS{45}Porque também não veio o Filho do homem para ser servido, mas para servir, e dar sua vida em resgate por muitos.
\VS{46}E vieram a Jericó. E saindo ele, e seus discípulos, e uma grande multidão de Jericó, estava Bartimeu, um cego, filho de Timeu, sentado junto ao caminho.
\VS{47}E ouvindo que era Jesus o nazareno, começou a clamar, e a dizer: Jesus, Filho de Davi! Tem misericórdia de mim!
\VS{48}E muitos o repreendiam, para que se calasse; mas ele clamava ainda mais: Filho de Davi! Tem misericórdia de mim!
\VS{49}Jesus parou, e disse que o chamassem. Então chamaram ao cego, dizendo-lhe: Tem bom ânimo, levanta-te, ele está te chamando.
\VS{50}Então ele lançou sua capa, levantou-se, e veio a Jesus.
\VS{51}Jesus, lhe perguntou: Que queres que eu te faça? E o cego lhe disse: Mestre, que eu veja.
\VS{52}Jesus lhe disse: Vai; a tua fé te salvou. E logo consegiu ver. Então seguiu Jesus pelo caminho.

\par }\Chap{11}{\PP \VerseOne{1}E quando se aproximaram de Jerusalém, de Betfagé, e de Betânia, junto ao monte das Oliveiras,
{\ADD{Jesus}} enviou dois de seus discípulos,
\VS{2}dizendo-lhes: Ide ao vilarejo que está adiante de vós; e assim que nela entrardes, achareis um jumentinho amarrado, sobre o qual ninguém se sentou; soltai-o, e trazei-o.
\VS{3}E se alguém vos disser: Por que fazeis isso?, dizei: O Senhor precisa dele, e logo o devolverá para cá.
\VS{4}Eles foram, e acharam o jumentinho amarrado à porta, do lado de fora em uma esquina, e o soltaram.
\VS{5}E alguns dos que ali estavam lhes perguntaram: Que fazeis, soltando o jumentinho?
\VS{6}Eles lhes disseram como Jesus
{\ADD{lhes}} havia mandado, e os deixaram ir.
\VS{7}Então trouxeram o jumentinho a Jesus. Lançaram sobre ele suas roupas, e
{\ADD{Jesus}} sentou-se sobre ele.
\VS{8}Muitos estendiam suas roupas pelo caminho, e outros cortavam ramos das árvores, e os espalhavam pelo caminho..
\VS{9}E os que iam adiante, e os que seguiam, clamavam: Hosana, bendito o que vem no Nome do Senhor!
\VS{10}Bendito o Reino que vem no Nome do Senhor,
{\ADD{o Reino}} do nosso pai Davi! Hosana nas alturas!
\VS{11}Jesus entrou em Jerusalém, e no Templo. E depois que ter visto tudo em redor, e sendo já tarde, ele saiu para Betânia com os doze.
\VS{12}E no dia seguinte, quando saíram de Betânia, ele teve fome.
\VS{13}E vendo de longe uma figueira que tinha folhas,
{\ADD{veio ver}} se acharia alguma coisa nela; mas ao chegar perto dela, nada achou, a não ser folhas, pois não era o tempo de figos.
\VS{14}Então Jesus lhe disse: Nunca mais ninguém coma fruto de ti! E seus discípulos ouviram isso.
\VS{15}Depois vieram a Jerusalém. E entrando Jesus no Templo, começou a expulsar os que vendiam e compravam no Templo; e revirou as mesas dos cambiadores, e as cadeiras dos que vendiam pombas.
\VS{16}E não consentia que ninguém levasse vaso algum pelo Templo.
\VS{17}E ensinava, dizendo-lhes: Não está escrito: Minha casa será chamada casa de oração de todas as nações?
{\RQ{Isaías 56:7
}} Mas vós fizestes dela esconderijo de ladrões!
\VS{18}Os escribas e os chefes dos sacerdotes ouviram isso, e buscavam uma maneira de o matar; pois o temiam, porque toda a multidão estava admirada do ensino dele.
\VS{19}E como já era tarde,
{\ADD{Jesus}} saiu fora da cidade.
\VS{20}E passando pela manhã, viram que a figueira estava seca desde as raízes.
\VS{21}Pedro se lembrou disso, e disse-lhe: Mestre, eis que a figueira que amaldiçoaste, se secou.
\VS{22}E respondendo Jesus, disse-lhes: Tende fé em Deus.
\VS{23}Porque em verdade vos digo que qualquer um que disser a este monte: Levanta-te, e lança-te no mar; e não duvidar em seu coração, mas crer que se fará o que diz, tudo o que disser lhe será feito.
\VS{24}Portanto eu vos digo que tudo o que pedirdes orando, crede que recebereis, e vós
{\ADD{o}} tereis.
\VS{25}E quando estiverdes orando, perdoai, se tendes algo contra alguém, para que o vosso Pai, que
{\ADD{está}} nos céus, vos perdoe vossas ofensas.
\VS{26}Mas se vós não perdoardes, também o vosso Pai, que
{\ADD{está}} nos céus, não vos perdoará vossas ofensas.
\VS{27}Depois voltaram a Jerusalém; e, enquanto ele andava pelo Templo, vieram a ele os chefes dos sacerdotes, os escribas, e os anciãos.
\VS{28}E disseram-lhe: Com que autoridade fazes estas coisas? E quem te deu esta autoridade, para fazerdes estas coisas?
\VS{29}Jesus lhes respondeu: Também eu vos farei uma pergunta, e respondei-me; então vos direi com que autoridade faço estas coisas.
\VS{30}O batismo de João era do céu ou dos homens? Respondei-me.
\VS{31}E eles argumentavam entre si, dizendo: Se dissermos do céu, ele dirá: Por que, pois, não crestes nele?
\VS{32}Porém, se dissermos dos homens, tememos ao povo, porque todos consideravam que João era verdadeiramente profeta.
\VS{33}Então responderam a Jesus: Não sabemos. E Jesus lhes replicou: Também eu não vos direi com que autoridade faço estas coisas.

\par }\Chap{12}{\PP \VerseOne{1}Então
{\ADD{Jesus}} começou a lhes dizer por parábolas: Um homem plantou uma vinha, cercou-a, fundou
{\ADD{nela}} uma prensa de uvas, edificou uma torre, e a arrendou a uns lavradores; e partiu-se daquela terra.
\VS{2}E, quando chegou o tempo, mandou um servo aos lavradores, para que recebesse dos lavradores do fruto da vinha.
\VS{3}Mas eles, tomando-o à força, feriram-no, e mandaram-no vazio.
\VS{4}E voltou a lhes mandar outro servo; e eles, apedrejando-o, feriram-no na cabeça e voltaram a mandá-lo maltratado.
\VS{5}E voltou a mandar outro, ao qual mataram, e
{\ADD{mandou}} muitos outros, e a uns feriram, e a outros mataram.
\VS{6}Tendo ele, pois, ainda um, o seu filho amado, mandou-lhes também por último a este, dizendo: Pelo menos respeitarão o meu filho.
\VS{7}Mas aqueles lavradores disseram entre si: Este é o herdeiro; vinde, e o matemos; então a herança será nossa.
\VS{8}E pegando dele, mataram{\ADD{-no}} , e lançaram{\ADD{-no}} fora da vinha.
\VS{9}Que, pois, o senhor da vinha fará? Ele virá, destruirá aos lavradores, e dará a vinha a outros.
\VS{10}Acaso ainda não lestes esta escritura: “A pedra que os edificadores rejeitaram, esta foi feita por cabeça de esquina.
\VS{11}Pelo Senhor foi feito isto, e é maravilhoso em nossos olhos”?
{\RQ{Salmos 118:22-23
}}
\VS{12}E buscavam prendê-lo, mas temiam a multidão; porque entendiam que era contra eles que dizia aquela parábola; então o deixaram e se retiraram.
\VS{13}E mandaram-lhe alguns dos fariseus e dos herodianos, para que o apanhassem em
{\ADD{alguma}} palavra.
\VS{14}Eles se aproximaram e lhe disseram: Mestre, sabemos que és homem de verdade, e não te interessa
{\ADD{agradar}} a ninguém, porque não te importas com a aparência humana, mas com verdade ensinas o caminho de Deus. É lícito pagar tributo a César, ou não? Devemos pagar, ou não devemos?
\VS{15}E ele, entendendo a hipocrisia deles, disse-lhes: Por que estais me tentando? Trazei-me uma moeda,
\FTNT{*}{{\FR 12:15 }
ou “denário”, moeda romana
} para que eu a veja.
\VS{16}E trouxeram. E perguntou-lhes: De quem é esta imagem, e a inscrição? E eles lhe disseram: De César.
\VS{17}Então Jesus, lhes respondeu: Dai, pois, a César o que é de César, e a Deus o que é de Deus. E admiram-se dele.
\VS{18}E vieram a ele os saduceus, que dizem não haver ressurreição, e perguntaram-lhe:
\VS{19}Mestre, Moisés nos escreveu que, se o irmão de alguém morresse, e deixasse mulher, mas não filhos, que seu irmão se casasse com a viúva
\FTNT{†}{{\FR 12:19 }
Lit. “a mulher dele” (isto é, do morto)
} , e gerasse descendência
\FTNT{‡}{{\FR 12:19 }
Lit. “levantasse semente”
} ao seu irmão.
\VS{20}Houve, pois, sete irmãos, e o primeiro casou-se com a mulher, e morrendo, não deixou descendente.
\VS{21}Casou-se com ela também o segundo, e morreu; e nem este deixou descendente; e o terceiro da mesma maneira.
\VS{22}Os sete casaram-se com ela, mas não deixaram descendente. Finalmente, depois de todos, morreu também a mulher.
\VS{23}Na ressurreição, pois, quando ressuscitarem, ela será a mulher de qual deles? Porque os sete a tiveram por mulher.
\VS{24}E Jesus lhes respondeu: Acaso não é por isso que errais, por não conhecerdes as Escrituras, nem o poder de Deus?
\VS{25}Pois, quando ressuscitarem dos mortos, nem se casarão, nem se darão em casamento; mas serão como os anjos que estão nos céus.
\VS{26}E quanto aos mortos que ressuscitarão, não lestes no livro de Moisés, como Deus lhe falou com a sarça, dizendo: Eu sou o Deus de Abraão, o Deus de Isaque, e o Deus de Jacó?
{\RQ{Êxodo 3:6
}}
\VS{27}Deus não é
{\ADD{Deus}} de mortos, mas de vivos. Portanto errais muito.
\VS{28}Então aproximou-se dele um dos escribas, que os havia ouvido discutir. Como ele sabia que
{\ADD{Jesus}} havia lhes respondido bem, perguntou-lhe: Qual é o primeiro mandamento de todos?
\VS{29}E Jesus lhe respondeu: O primeiro de todos os mandamentos é: “Ouve Israel, o Senhor nosso Deus é o único Senhor;
\VS{30}amarás, pois, ao Senhor teu Deus com todo o teu coração, com toda a tua alma, com todo o teu entendimento, e com todas as tuas forças.”
{\RQ{Deuteronômio 6:4-5
}} Este é o primeiro mandamento.
\VS{31}E o segundo é semelhante: “Amarás ao teu próximo como a ti mesmo”.
{\RQ{Levítico 19:18
}} Não há outro mandamento maior que estes.
\VS{32}E o escriba lhe disse: Muito bem, Mestre, com verdade disseste que há um só Deus, e não há outro além dele.
\VS{33}E que amá-lo com todo o coração, com todo o entendimento, com toda a alma, e com todas as forças; e amar ao próximo como a si mesmo, é mais de que todos os holocaustos e sacrifícios.
\VS{34}Jesus viu que ele havia respondido sabiamente, e disse-lhe: Não estás longe do Reino de Deus. E ninguém mais ousou lhe perguntar.
\VS{35}E Jesus respondia e dizia, enquanto ensinava no templo: Como os escribas dizem que o Cristo é filho de Davi?
\VS{36}Porque o mesmo Davi disse pelo Espírito Santo: Disse o Senhor ao meu Senhor, senta-te à minha direita, até que eu ponha os teus inimigos debaixo dos teus pés.
{\RQ{Salmos 110:1
}}
\FTNT{§}{{\FR 12:36 }
Lit. por escabelo (um tipo de banquinho) dos teus pés
}
\VS{37}Ora,
{\ADD{se}} o próprio Davi o chama de Senhor, como, pois, é so eu filho?E a grande multidão o ouvia de boa vontade.
\VS{38}E dizia-lhes em seu ensino: Cuidado com os escribas, que gostam de andar com roupas compridas, das saudações nas praças,
\VS{39}das primeiras cadeiras nas sinagogas, e dos primeiros assentos nas ceias.
\VS{40}que devoram as casas das viúvas, com pretexto de longas orações. Esses receberão mais grave condenação.
\VS{41}Quando Jesus estava sentado de frente à arca do tesouro, observava como a multidão lançava dinheiro na arca do tesouro; e muitos ricos lançavam muito.
\VS{42}E uma pobre viúva veio, e lançou duas pequenas moedas, de pouco valor.
\FTNT{*}{{\FR 12:42 }
duas pequenas moedas, de pouco valor
Lit. dois leptos, que valiam um quadrante
}
\VS{43}Então
{\ADD{Jesus}} chamou a si os seus discípulos, e disse-lhes: Em verdade vos digo que esta pobre viúva lançou mais que todos os que lançaram na arca do tesouro;
\VS{44}Porque todos lançaram daquilo que lhes sobra; mas esta, da sua pobreza lançou tudo o que tinha, todo o seu sustento.

\par }\Chap{13}{\PP \VerseOne{1}Quando ele estava saindo do templo, um de seus discípulos lhe disse: Mestre, olha que pedras, e que edifícios!
\VS{2}Jesus respondeu, dizendo-lhe: Vês estes grandes edifícios? Não será deixada pedra sobre pedra, que não seja derrubada.
\VS{3}Depois, sentando-se ele no monte das Oliveiras, de frente ao templo, Pedro, Tiago, João, e André perguntaram-lhe à parte:
\VS{4}Dize-nos, quando serão essas coisas? E que sinal haverá de quando todas essas coisas se cumprirão?
\VS{5}E Jesus, respondendo-lhes, começou a dizer: Cuidado! ninguém vos engane;
\VS{6}porque muitos virão em meu nome, dizendo: “Sou eu”, e enganarão a muitos.
\VS{7}E quando ouvirdes de guerras e de rumores de guerras, não vos perturbeis; porque
{\ADD{assim}} deve acontecer; mas ainda não será o fim.
\VS{8}Porque nação se levantará contra nação, e reino contra reino, e haverá terremotos em muitos lugares, e haverá fomes e tumultos. Estes serão os princípios das dores.
\VS{9}Mas cuidai de vós mesmos; porque vos entregarão em tribunais e em sinagogas; sereis açoitados, e sereis levados diante de governadores e reis, por causa de mim, para que lhes haja testemunho.
\VS{10}Mas antes o Evangelho deve ser pregado entre todas as nações.
\VS{11}Porém, quando vos levarem para vos entregar, não estejais ansiosos antecipadamente do que deveis dizer, nem o planejeis; mas o que naquela hora for dado, isso falai. Porque não sois vós os que falais, mas o Espírito Santo.
\VS{12}E o irmão entregará ao irmão à morte, e o pai ao filho; e os filhos se levantarão contra os pais, e os matarão.
\VS{13}E sereis odiados por todos por causa do meu nome; mas quem perseverar até o fim, esse será salvo.
\VS{14}E quando virdes a abominação da desolação, que foi dita pelo profeta Daniel, estar onde não deve, (quem lê, entenda), então os que estiverem na Judeia, fujam para os montes.
\VS{15}E quem estiver sobre telhado, não desça para a casa, nem entre para tomar alguma coisa de sua casa.
\VS{16}E quem estiver no campo, não volte atrás, para tomar sua roupa.
\VS{17}Mas ai das grávidas, e das que amamentarem naqueles dias!
\VS{18}Orai, porém, para que a vossa fuga não aconteça no inverno.
\VS{19}Porque aqueles dias serão de tal aflição, qual nunca foi desde o princípio da criação das coisas, que Deus criou, até agora, nem jamais será.
\VS{20}E se o Senhor não encurtasse aqueles dias, ninguém
\FTNT{*}{{\FR 13:20 }
Lit. “nenhuma carne”
} se salvaria; mas por causa dos escolhidos, que escolheu, ele encurtou aqueles dias.
\VS{21}Então se alguém vos disser: “Eis aqui o Cristo”; ou “ei-lo ali”, não creiais nele.
\VS{22}Porque se levantarão falsos cristos, e falsos profetas, que farão sinais e prodígios, para enganar, se possível, até os escolhidos.
\VS{23}Vós, porém, tende cuidado; eis que tenho vos dito tudo com antecedência.
\VS{24}Mas naqueles dias, depois daquela aflição, o sol se escurecerá, e a lua não dará seu brilho;
\VS{25}as estrelas do céu cairão, e as forças que estão nos céus se abalarão.
\VS{26}Então verão o Filho do homem vir nas nuvens, com grande poder e glória.
\VS{27}Ele então enviará os seus anjos, e ajuntará os seus escolhidos dos quatro ventos, desde a extremidade da terra, até a extremidade do céu.
\VS{28}E aprendei a parábola da figueira: Quando o seu ramo já vai ficando tenro, e brota folhas, bem sabeis que o verão já está perto.
\VS{29}Assim também vós, quando virdes suceder estas coisas, sabei que já está perto, às portas.
\VS{30}Em verdade vos digo que não passará esta geração, até que todas estas coisas aconteçam.
\VS{31}O céu e a terra passarão, mas as minhas palavras em maneira nenhuma passarão.
\VS{32}Porém daquele dia e hora ninguém sabe, nem os anjos que estão no céu, nem o Filho, a não ser somente o Pai.
\VS{33}Olhai, vigiai, e orai; porque não sabeis quando será o tempo.
\VS{34}Assim como o homem que, partindo-se para fora de sua terra, deixou sua casa, e deu autoridade aos seus servos, e a cada um o seu trabalho, e mandou ao porteiro que vigiasse,
\VS{35}assim também vigiai, porque não sabeis quando virá o Senhor da casa; se à tarde, se à meia-noite, se ao cantar do galo, se pela manhã,
\VS{36}para que ele não venha de repente, e vos ache dormindo.
\VS{37}E as coisas que vos digo, digo a todos: Vigiai.

\par }\Chap{14}{\PP \VerseOne{1}E dali a dois dias era a Páscoa, e
{\ADD{a festa}} dos pães sem fermento; e os chefes dos sacerdotes, e os escribas buscavam um meio de prendê-lo através de um engano, e o matarem.
\VS{2}Diziam, porém: “Não na festa, para que não aconteça tumulto entre o povo.”
\VS{3}E estando ele em Betânia, na casa de Simão o Leproso, sentado
{\ADD{à mesa}} , veio uma mulher, que tinha um vaso de alabastro, de óleo perfumado de nardo puro, de muito preço, e quebrando o vaso de alabastro, derramou-o sobre a cabeça dele.
\VS{4}E houve alguns que em si mesmos se indignaram, e disseram: “Para que foi feito este desperdício do óleo perfumado?
\VS{5}Porque isto podia ter sido vendido por mais de trezentos denários, e seria dado aos pobres.” E reclamavam contra ela.
\VS{6}Porém Jesus disse: “Deixai-a; por que a incomodais? Ela me fez boa obra.
\VS{7}Porque pobres sempre tendes convosco; e quando quiserdes, podeis lhes fazer bem; porém a mim, nem sempre me tendes.
\VS{8}Esta fez o que podia; adiantou-se para ungir o meu corpo, para a sepultura.
\VS{9}Em verdade vos digo, que onde quer que em todo o mundo este Evangelho for pregado, também o que esta fez será dito em sua memória.”
\VS{10}E Judas Iscariotes, um dos doze, foi aos chefes dos sacerdotes, para o entregar a eles.
\VS{11}E eles ouvindo, alegraram-se; e prometeram lhe dar dinheiro; e buscava como o entregaria em tempo oportuno.
\VS{12}E o primeiro dia dos pães sem fermento, quando sacrificavam
{\ADD{o cordeiro da}} Páscoa, seus discípulos lhe disseram: “Onde queres que vamos preparar para comerdes a Páscoa?”
\VS{13}E mandou dois de seus discípulos, e disse-lhes: “Ide à cidade, e um homem que leva um cântaro de água vos encontrará, a ele segui.
\VS{14}E onde quer que ele entrar, dizei ao senhor da casa: O Mestre diz: Onde está o cômodo onde comerei Páscoa com meus discípulos?
\VS{15}E ele vos mostrará um grande salão, ornado e preparado; ali fazei os preparativos para nós.”
\VS{16}E seusdiscípulos saíram, e vieram à cidade, e acharam como havia lhes dito, e prepararam a Páscoa.
\VS{17}E chegada a tarde, veio com os doze.
\VS{18}E quando se sentaram
{\ADD{à mesa}} , e comeram, Jesus disse: “Em verdade vos digo, que um de vós, que está comendo comigo, me trairá.”
\VS{19}E eles começaram a se entristecer, e a lhe dizer um após outro: “Por acaso sou eu?” E outro: “Por acaso sou eu?”
\VS{20}Porém ele lhes respondeu: “{\ADD{É}} um dos doze, o que está molhando
{\ADD{a mão}} comigo no prato.
\VS{21}Em verdade o Filho do homem vai, como está escrito sobre ele; mas ai daquele homem por quem o Filho do homem é traído; bom lhe fosse ao tal homem não haver nascido.”
\VS{22}E enquanto eles comiam, Jesus tomou o pão; e bendizendo, partiu-o, deu-lhes, e disse: “Tomai, comei, isto é o meu corpo.”
\VS{23}E tomando o copo, e dando graças, deu-lhes; e todos beberam dele.
\VS{24}E disse-lhes: “Isto é o meu sangue, o do novo testamento, que é derramado por muitos.
\VS{25}Em verdade vos digo, que não beberei mais do fruto da vide, até aquele dia, quando o beber novo no Reino de Deus.”
\VS{26}E depois de cantarem um hino, saíram para o monte das Oliveiras.
\VS{27}E Jesus lhes disse: “Todos vós falhareis
\FTNT{*}{{\FR 14:27 }
Ou: me abandonareis.
} comigo esta noite; porque está escrito: Ferirei ao pastor, e as ovelhas serão dispersas.
{\RQ{Zacarias 13:7
}}
\VS{28}Mas depois de eu haver ressuscitado, irei adiante de vós para a Galileia.”
\VS{29}E Pedro lhe disse: “Ainda que todos falhem, eu não
{\ADD{falharei}} .”
\VS{30}Jesus lhe disse: “Em verdade te digo, que hoje, nesta noite, antes que o galo cante duas vezes, me negarás três vezes.”
\VS{31}Mas
{\ADD{Pedro}} , insistindo, dizia: “Ainda que me seja necessário morrer contigo, em maneira nenhuma te negarei.” E todos diziam também da mesma maneira.
\VS{32}E vieram ao lugar, cujo nome era Getsêmani, e disse a seus discípulos: “Sentai-vos aqui, enquanto eu oro.”
\VS{33}E tomou consigo Pedro, Tiago, e João, e começou a ficar muito apavorado e angustiado.
\VS{34}E disse-lhes: “Minha alma totalmente está triste até a morte; ficai-vos aqui, e vigiai.”
\VS{35}Então ele foi um pouco mais adiante, e prostrou-se em terra. E orou, que se fosse possível, afastasse dele aquela hora.
\VS{36}E disse: “Aba, Pai, todas as coisas te são possíveis; passa de mim este copo; porém não
{\ADD{se faça}} o que eu quero, mas sim o que tu
{\ADD{queres}} .”
\VS{37}Depois veio de volta, e os achou dormindo; e disse a Pedro: “Simão, estás dormindo? Não podes vigiar uma hora?
\VS{38}Vigiai, e orai, para que não entreis em tentação; o espírito em verdade
{\ADD{está}} pronto, mas a carne
{\ADD{é}} fraca.”
\VS{39}E depois que se foi novamente, orou, dizendo as mesmas palavras.
\VS{40}Quando voltou, achou-os outra vez dormindo; porque os olhos deles estavam pesados, e não sabiam o que lhe responder.
\VS{41}E veio a terceira vez, e disse-lhes: “Ainda estais dormindo e descansando?
\FTNT{†}{{\FR 14:41 }
Trad. alt. Agora dormi e descansai.
} Basta, chegada é a hora. Eis que o Filho do homem é entregue em mãos dos pecadores.
\VS{42}Levantai-vos, vamos; eis que o que me trai está perto.
\VS{43}E logo, enquanto ele ainda estava falando, veio Judas, que era um dos doze, e com ele uma grande multidão, com espadas e bastões, da parte dos chefes dos sacerdotes, dos escribas, dos anciãos.
\VS{44}E o que o traía lhes tinha dado um sinal comum, dizendo: “Ao que eu beijar, é esse; prendei-o, e levai-o em segurança.”
\VS{45}E quando veio, logo foi-se a ele, e disse-lhe: “Rabi, Rabi”, e o beijou.
\VS{46}Então o agarraram,
\FTNT{‡}{{\FR 14:46 }
Lit. lançaram suas mãos nele
} e o prenderam.
\VS{47}E um dos que estavam ali presentes puxou a espada, feriu o servo do sumo sacerdote, e cortou-lhe a orelha.
\VS{48}Jesus começou a lhes dizer: “Viestes me prender com espadas e bastões, como
{\ADD{se eu fosse}} um bandido?
\VS{49}Todo dia eu estava convosco no templo, ensinando, e não me prendestes; mas
{\ADD{assim se faz}} para que as Escrituras se cumpram.”
\VS{50}Então todos o deixaram, e fugiram.
\VS{51}E certo rapaz o seguia, envolto num lençol sobre o
{\ADD{corpo}} nu. E os rapazes o agarraram.
\VS{52}Mas ele largou o lençol, e fugiu deles nu.
\VS{53}E levaram Jesus ao sumo sacerdote; e ajuntaram-se a ele todos os chefes dos sacerdotes, os anciãos, e os escribas.
\VS{54}E Pedro o seguiu de longe até dentro do pátio do sumo sacerdote, e ficou sentado com os oficiais, esquentando-se ao fogo.
\VS{55}E os chefes dos sacerdotes, e todo o supremo conselho
\FTNT{§}{{\FR 14:55 }
supremo conselho
lit. sinédrio – o mais importante conselho ou tribunal para os judeus
} buscavam
{\ADD{algum}} testemunho contra Jesus, para o matarem, e não
{\ADD{o}} achavam.
\VS{56}Porque muitos davam falso testemunho contra ele, mas os testemunhos não concordavam entre si.
\VS{57}E alguns se levantavam e davam falso testemunho contra ele, dizendo:
\VS{58}“Nós o ouvimos dizer: Eu derrubarei este templo feito por mãos, e em três dias construirei outro feito não por mãos.
\VS{59}E nem assim o testemunho deles era concordante.
\VS{60}Então o sumo sacerdote levantou-se no meio, e perguntou a Jesus: “Não respondes nada? O que estes testemunham contra ti?”
\VS{61}Mas ele ficou calado, e nada respondeu. O sumo sacerdote voltou a lhe perguntar: “És tu o Cristo, o Filho daquele que é Bendito?”
\VS{62}Jesus respondeu: “Eu sou. E vereis o Filho do homem sentado à direita do Poderoso
\FTNT{*}{{\FR 14:62 }
Lit. Poder
} , e vindo com as nuvens do céu.”
\VS{63}E o sumo sacerdote, rasgando suas roupas, disse: “Para que mais necessitamos de testemunhas?
\VS{64}Vós ouvistes a blasfêmia. Que vos parece?” E todos o condenaram como culpado de morte.
\VS{65}E alguns começaram a cuspir nele, e a cobrir o rosto dele; e a dar-lhe de socos, e dizer-lhe: “Profetiza”. E os oficiais lhe davam bofetadas.
\VS{66}E, enquanto Pedro estava no pátio abaixo, veio uma das servas do sumo sacerdote.
\VS{67}Quando ela viu Pedro, que estava sentado esquentando-se, olhou para ele, e disse: “Também tu estavas com Jesus, o Nazareno.”
\VS{68}Mas ele negou, dizendo: “Não
{\ADD{o}} conheço, nem sei o que dizes.”; e saiu para o alpendre. Então o galo cantou.
\VS{69}A serva o viu outra vez, e começou a dizer aos que ali estavam: “Este é um deles”.
\VS{70}Mas ele o negou de novo. E pouco depois, outra vez os que ali estavam disseram a Pedro: “Verdadeiramente tu és um deles; pois também és galileu”, e a tua fala é semelhante”.
\VS{71}Então ele começou a amaldiçoar e a jurar,
{\ADD{dizendo}} : “Não conheço esse homem de quem dizeis.”
\VS{72}O galo cantou a segunda vez. E Pedro se lembrou da palavra que Jesus havia lhe dito: “Antes que o galo cante duas vezes, tu me negarás três vezes.” Então ele retirou-se dali e chorou.

\par }\Chap{15}{\PP \VerseOne{1}E logo ao amanhecer, os chefes dos sacerdotes tiveram uma reunião com os anciãos, com os escribas, e com todo o supremo conselho;
\FTNT{*}{{\FR 15:1 }
supremo conselho
lit. sinédrio – o mais importante conselho ou tribunal para os judeus
} e, amarrando Jesus, levaram
{\ADD{-no}} e
{\ADD{o}} entregaram a Pilatos.
\VS{2}E Pilatos lhe perguntou: És tu o Rei dos Judeus? E ele lhe respondeu: Tu o dizes.
\VS{3}E os chefes dos sacerdotes o acusavam de muitas coisas, porém ele nada respondia.
\VS{4}E outra vez Pilatos lhe perguntou: Não respondes nada? Olha quantas
{\ADD{coisas}} testemunham contra ti!
\VS{5}Mas Jesus nada mais respondeu, de maneira que Pilatos ficou surpreso.
\VS{6}Na festa,
{\ADD{Pilatos}} lhes soltava um preso, qualquer um que pedissem.
\VS{7}E havia um chamado Barrabás, preso com
{\ADD{outros}} revoltosos, que numa rebelião havia cometido uma morte.
\VS{8}E a multidão, gritando, começou a pedir, como sempre lhes havia feito.
\VS{9}E Pilatos lhes respondeu: Quereis que eu vos solte o Rei dos Judeus?
\VS{10}(Porque ele sabia que os chefes dos sacerdotes haviam o entregue por inveja).
\VS{11}Mas os chefes dos sacerdotes incitaram a multidão, para que, em vez disso, lhes soltasse Barrabás.
\VS{12}E Pilatos, respondendo, disse-lhes outra vez: Que, pois, quereis que eu faça com aquele a quem chamais Rei dos Judeus?
\VS{13}E eles voltaram a clamar: Crucifica-o!
\VS{14}Mas Pilatos lhes disse: Por quê? Que mal ele fez? E eles gritavam ainda mais: Crucifica-o!
\VS{15}Então Pilatos, querendo satisfazer à multidão, soltou-lhes Barrabás; e entregou Jesus açoitado, para que fosse crucificado.
\VS{16}E os soldados o levaram para o pátio, que é o pretório; e convocaram toda a tropa.
\VS{17}E o vestiram de púrpura; teceram uma coroa de espinhos, e puseram nele.
\VS{18}E começaram a saudá-lo: Viva! Ó Rei dos Judeus!
\VS{19}E feriram a sua cabeça com uma cana, cuspiram nele, e, ajoelhados, o adoraram.
\VS{20}Quando o escarneceram, despiram-lhe a púrpura, vestiram-no com as suas próprias roupas, e o levaram afora, para o crucificarem.
\VS{21}E forçaram um Simão cireneu, que estava passando, vindo do campo, o pai de Alexandre e de Rufo, para que levasse sua cruz.
\VS{22}E o levaram ao lugar de Gólgota, que traduzido é: o lugar da caveira.
\VS{23}E ofereceram-lhe a beber vinho misturado com mirra; mas ele não o tomou.
\VS{24}E havendo o crucificado, repartiram a roupas dele, lançando-lhes sortes, para o que cada um levaria.
\VS{25}Era a hora terceira, e o crucificaram.
\VS{26}E a descrição de sua acusação estava acima escrita: O REI DOS JUDEUS.
\VS{27}E crucificaram com ele dois ladrões,
\FTNT{†}{{\FR 15:27 }
Ou: rebeldes
} um à sua direita, e outro à esquerda.
\VS{28}E cumpriu-se a Escritura que diz: E foi contado com os malfeitores.
\VS{29}E os que passavam blasfemavam dele, balançando suas cabeças, e dizendo: Ah! tu que derrubas o templo, e em três dias o edificas,
\VS{30}salva a ti mesmo, e desce da cruz!
\VS{31}E da mesma maneira também os chefes dos sacerdotes, com os escribas, diziam aos outros, escarnecendo: Ele salvou a outros, a si mesmo não pode salvar!
\VS{32}Que o Cristo, o Rei de Israel, desça agora da cruz, para que o vejamos, e creiamos! Os que estavam crucificados com ele também o insultavam.
\VS{33}E vinda a hora sexta, vieram trevas sobre toda a terra, até a hora nona.
\VS{34}E na hora nona, Jesus exclamou em alta voz: ELOÍ, ELOÍ, LAMÁ SABACTÂNI, que traduzido é: Deus meu, Deus meu, por que me desamparaste?
\VS{35}E alguns dos que ali estavam, quando ouviram, disseram: Eis que ele está chamando Elias.
\VS{36}E um correu, encheu de vinagre uma esponja; e pondo-a em uma cana, dava-lhe de beber, dizendo: Deixai, vejamos se Elias virá tirá-lo.
\VS{37}E Jesus bradou em grande voz, então expirou.
\VS{38}E o véu do Templo se rasgou em dois do alto abaixo.
\VS{39}E o centurião que estava ali diante dele, vendo que havia expirado bradando assim, disse: Verdadeiramente este homem era Filho de Deus.
\VS{40}E também estavam ali
{\ADD{algumas}} mulheres olhando de longe, entre as quais estava também Maria Madalena, e Maria (mãe de Tiago o menor e de José), e Salomé;
\VS{41}as quais também, quando ele estava na Galileia, o seguiam, e o serviam; e outras muitas, que haviam subido com ele a Jerusalém.
\VS{42}E quando já vinha o final da tarde, porque era a preparação, que é o dia antes de sábado,
\VS{43}Veio José de Arimateia, honrado membro do conselho, que também esperava o Reino de Deus, e com ousadia foi a Pilatos, e pediu o corpo de Jesus.
\VS{44}Pilatos se surpreendeu de que já fosse morto. E chamou a si o centurião, e perguntou-lhe se já era morto já havia muito tempo.
\VS{45}Quando ele recebeu a explicação do centurião, deu o corpo a José,
\VS{46}o qual comprou um lençol fino, e tirando-o
{\ADD{da cruz}} , envolveu-o no lençol fino. Em seguida, ele o pôs num sepulcro escavado em uma rocha, e rolouuma pedra à porta do sepulcro.
\VS{47}Maria Madalena e Maria
{\ADD{mãe}} de José olharam onde o puseram.

\par }\Chap{16}{\PP \VerseOne{1}Havendo passado o sábado, Maria Madalena, e Maria
{\ADD{mãe}} de Tiago, e Salomé, compraram especiarias, para irem o ungir.
\VS{2}E de manhã muito cedo, o primeiro dia da semana, vieram ao sepulcro, ao nascer do sol.
\VS{3}E diziam umas às outras: Quem rolará para nós a pedra da porta do sepulcro?
\VS{4}Pois era muito grande. E quando observaram, viram que já a pedra estava havia sido rolada;
\VS{5}e entrando no sepulcro, viram um rapaz sentado à direita, vestido com uma roupa comprida branca; e elas se espantaram.
\VS{6}Mas ele lhes disse: Não vos espanteis. Vós buscais a Jesus Nazareno crucificado. Ele já ressuscitou, não está aqui. Eis aqui o lugar onde o puseram.
\VS{7}Porém ide, dizei a seus discípulos e a Pedro, que ele vai adiante de para a Galileia; ali o vereis, como ele vos disse.
\VS{8}Então elas, saindo apressadamente, fugiram do sepulcro, porque temor e espanto as haviam tomado; e não disseram nada a ninguém, porque tinham medo.
\VS{9}E havendo ressuscitado pela manhã, no primeiro
{\ADD{dia}} da semana,
{\ADD{Jesus}} apareceu primeiramente a Maria Madalena, da qual havia expulsado sete demônios.
\VS{10}Ela se foi, e anunciou aos que estiveram com ele, que estavam tristes e chorando.
\VS{11}Quando ouviram que ele estava vivo, e que havia sido visto por ela, eles não creram.
\VS{12}E depois ele apareceu em outra forma a dois deles, que estavam indo pelo caminho ao campo.
\VS{13}E estes foram anunciar aos outros; mas também não creram neles.
\VS{14}Finalmente ele apareceu aos onze, estando eles sentados juntos; e repreendeu pela incredulidade e dureza de coração deles, por não terem crido nos que já o haviam visto ressuscitado.
\VS{15}E disse-lhes: Ide por todo o mundo, pregai o Evangelho a toda criatura.
\VS{16}Quem crer e for batizado será salvo; mas quem não crer será condenado.
\VS{17}E estes sinais seguirão aos que crerem: em meu nome expulsarão demônios; falarão novas línguas;
\VS{18}pegarão serpentes com as mãos; e se beberem algo mortífero, não lhes fará dano algum; porão as mãos sobre os enfermos, e sararão.
\VS{19}Então o Senhor, depois de haver lhes falado, foi recebido acima no céu, e sentou-se à direita de Deus.
\VS{20}E eles saíram e pregaram por todas as partes,
{\ADD{com}} o Senhor operando com eles, e confirmando a palavra com os sinais que se seguiam. Amém.\par }